\section{Related Work}\label{sec:related}

\paragraph{Guqin notation and performance-oriented symbolic music.}
Recent work on jianzipu has focused on recovering structured notation from glyphs and manuscript pages \cite{li2026zeroshotjianzi,li2026jztab}. Earlier language-model-based work also studied transcription from jianzipu into modern notation \cite{shi2024guqintranscription}. Our setting reverses this direction: starting from melodic events, the system must choose how to realize them on the instrument. A related problem arises in guitar tablature inference, where a pitch may admit several string--fret assignments and playability is assessed by musicians \cite{edwards2024miditotab}. Work on piano fingering likewise models hand position and execution difficulty \cite{iman2024pianofingering}. At the composition level, YuE2 generates a readable melody-and-chord score before synthesizing full-song audio; the score supports direct editing and revisions driven by external language models in response to user feedback \cite{yuan2026yue2}. \system studies how a given score is realized through guqin performance actions, including state-dependent gestures and compact notation whose interpretation depends on tuning and prior actions.

\paragraph{Tool-interactive planning and revision.}
ReAct interleaves reasoning and environment actions, while Toolformer and ToolLLM investigate tool-use learning and data construction \cite{yao2023react,schick2023toolformer,qin2024toolllm}. Self-Refine and Reflexion use iterative feedback to improve model behavior; CRITIC specifically studies revision driven by external tool outputs \cite{madaan2023selfrefine,shinn2023reflexion,gou2024critic}. \system uses a shared executable notation state to check pitch realizations and issue conservative, event-level warnings. These checks inform subsequent edits but do not determine musical suitability on their own.

\paragraph{Learning from interaction trajectories.}
AgentTuning and FireAct demonstrate the value of supervised interaction trajectories for language agents \cite{zeng2024agenttuning,chen2023fireact}. More closely related, RetroGen reconstructs process supervision from final expert artifacts for evidence-grounded generation \cite{huang2026retrogen}. Our teacher likewise has access to a target artifact during offline construction, but is required to execute edits in the student-facing guqin runtime; the retained training context contains only public actions and observations. GEPA instead optimizes prompts by reflecting on execution feedback without updating model weights, providing a distinct comparison to trajectory fine-tuning \cite{agrawal2026gepa}.
